%======================================================================
\section*{(g) Added after the deadline extension}
%======================================================================
The lab session showed that two camera settings decide whether a noise measurement is right: with the camera's default
8-bit ADC a Mono12 frame has only about 240 levels (steps of 16\,DN, larger than the read noise), and with the default
black level the dark level sits at 0\,DN and half of the dark noise is clipped. The connector fixed both at start-up.
They are now \textbf{tools the agent can use}, together with two more, so the MCP server has 13 tools (table in (b)):

\begin{center}\small
\begin{tabular}{@{}p{0.2\linewidth}p{0.26\linewidth}p{0.27\linewidth}p{0.2\linewidth}@{}}
\toprule
tool & what it does & lab camera (GenICam node) & simulated camera\\ \midrule
\texttt{set\_black\_level} & offset added before digitisation; returns the applied value &
  \texttt{BlackLevel} (integer; a larger value gives a lower dark level: 240 $\to$ 0\,DN, 180 $\to$ 62\,DN) &
  the dark level in DN\\
\texttt{set\_adc\_bit\_depth} & 8 or 12 bits of the ADC & \texttt{ADCBitDepth} (\texttt{Bits\_8}/\texttt{Bits\_12}),
  changed with the acquisition stopped & rounds to 256 levels of the 12-bit scale\\
\texttt{read\_temperature} & device temperature in $^\circ$C (dark current depends on it) & \texttt{DeviceTemperature} &
  warms from 38 to 41\,$^\circ$C in about 10 minutes\\
\texttt{save\_tiff} & a capture as a 16-bit multi-page TIFF for ImageJ/Fiji & on the server (Pillow) & the same\\
\bottomrule
\end{tabular}
\end{center}

The new settings are part of the camera interface (\texttt{camera/base.py}), so every connector offers them or refuses
them with a reason. A restarted server restores them from \texttt{calls.jsonl} like exposure and gain, every capture
reports them in its settings, and in the hardware turn \texttt{read\_temperature} and \texttt{save\_tiff} are commands
after which no frame is captured; the VERIFIED record of a hardware turn lists the temperatures read and the TIFF files.

\textbf{Packaging.} \texttt{hw3.py} is one entry point: \texttt{test} (all offline tests with a summary),
\texttt{check-camera} (the MVS SDK, then the plugged-in camera), \texttt{tools} (the server's tools as an MCP client sees
them), \texttt{run}, \texttt{live}, \texttt{demo}, and \texttt{mcp-config} (the entry that adds the camera server to any
MCP client, e.g.\ Claude Desktop). Every model call through the Claude Code CLI is now counted (calls, input and output
tokens, the CLI's cost estimate), and \texttt{model@low} lowers the model's thinking effort (36\,\% fewer output
tokens in a test). The MCP client always uses the Windows proactor event loop, so the camera server also starts inside
other programs, such as the web app of Homeworks~4/5.

\textbf{Tests} (\texttt{python hw3.py test}, simulated camera; all three suites pass). With the 8-bit ADC every value of a
dark frame is a multiple of $4095/255$ (16, 32, 48, 64\,DN, \dots); a black level of 100 gives a dark level of
100.0\,DN; the TIFF has one page per frame, equal to the frames; the device temperature reads 38.07\,$^\circ$C; and a
restarted server restores the 8-bit ADC, the black level of 100 and the gain of 6\,dB.

\textbf{Used by the agent.} In a run of the lab assistant of Homeworks~4/5 on the simulated camera (request: measure the
noise model and predict the focusing enhancement with it), the instrument team itself set the ADC to 12 bits and the
black level to 40 before its photon-transfer measurement, switched the light off with \texttt{set\_illumination} for 50
dark frames, and reported a conversion gain of 2.401\,e$^-$/DN (true value 2.40), a dark noise of 2.52\,DN and a full well
of at least 9737\,e$^-$ (true 10\,000\,e$^-$, not reached up to 30\,ms). It marked its read noise of 6.04\,e$^-$ (true
6.0\,e$^-$) as not verified because a cross-check failed; the optics team then predicted an enhancement of about 5.4 to
6.4 at 64 segments and 1000 iterations with this camera model.

\textbf{On the lab camera} the new connector functions use the same SDK calls as the connector's start-up, which ran at
the lab (\texttt{SetIntValueEx} for \texttt{BlackLevel}, \texttt{SetEnumValue} for \texttt{ADCBitDepth}), plus
\texttt{GetFloatValue} for \texttt{DeviceTemperature}; \texttt{tests/check\_real\_camera.py} now also checks them (black
level up and back, 8- and 12-bit frames, temperature). This check runs at the next lab session (18 October), before the
live demonstration of Homeworks~4/5.
